\providecommand{\atlasref}[2]{#2}
\section*{Literatura i status dowodów}
\addcontentsline{toc}{section}{Literatura i status dowodów}
\begin{enumerate}[label={[\arabic*]},leftmargin=*]
\item S. Saks, \emph{Theory of the Integral}, 2. wyd., Warszawa 1937; rozdziały o całkach Denjoya i Perrona.
\item R. A. Gordon, \emph{The Integrals of Lebesgue, Denjoy, Perron, and Henstock}, Graduate Studies in Mathematics 4, American Mathematical Society, 1994.
\item A. Fonda, \emph{The Kurzweil--Henstock Integral for Undergraduates}, Birkh\"{a}user, 2018.
\item M. Tvrd\'{y}, \emph{Kurzweil--Stieltjes Integration}, \url{https://users.math.cas.cz/~tvrdy/ks-dekk.pdf}.
\item Y. Kubota, \emph{A constructive definition of the approximately continuous Denjoy integral}, Canad. Math. Bull. 15 (1972), 103--109.
\item Materiały Biblioteki matematycznej: Analiza matematyczna I (rozdziały 19, 29, 36), Teoria miary (rozdziały 06--08, 14, 17--18, 20--21, 23), Analiza funkcjonalna (rozdział 12). Wykorzystano opracowane tam dowody w odpowiednich częściach niniejszego skryptu.
\end{enumerate}

\paragraph*{Zakres samodzielności dowodów.}
Bezpośrednio wykazano twierdzenia o sumach całkowych, wybrane kryteria, rachunki z funkcjami skokowymi, lemat uzupełniania podziałów, proste twierdzenia Perrona oraz własności wynikające z funkcji pierwotnych. Charakteryzację wszystkich pierwotnych KH przez $\AC G_*$ udowodniono w twierdzeniu \atlasref{AM1-38::it:thm:kh-acgstar-characterization}{charakteryzacja funkcji pierwotnych KH}, z wykorzystaniem wskazanych wyników teorii miary i całki Lebesgue'a. Twierdzenie Riesza--Markowa--Kakutaniego (twierdzenie \atlasref{ME-21::thm:riesz-positive-21}{ME-21: twierdzenie Riesza--Markowa--Kakutaniego}) udowodniono dla dowolnych lokalnie zwartych przestrzeni Hausdorffa przez zawartość zbiorów otwartych, miarę zewnętrzną i kryterium Carathéodory'ego; zachowano również wcześniejszy dowód metryczny. Dowody istnienia i jednoznaczności miary Haara bez założeń przeliczalności znajdują się w \atlasref{ME-23::thm:haar-main-23}{ME-23: twierdzenie Haara}. Wykorzystują wcześniejsze wyniki topologii i teorii miary, w szczególności twierdzenie Carathéodory'ego, przywoływane w tekście. Równoważność całek Perrona i HK udowodniono w twierdzeniu \atlasref{AM1-39::it:thm:perron-hk}{równoważność całek Perrona i KH}, z wykorzystaniem małej ciągłej korekty na zbiorze wyjątkowym (lemat \atlasref{AM1-39::it:lem:perron-null-correction}{korekta na zbiorze wyjątkowym}). Istnienie i jednoznaczność miary w skończonej wersji twierdzenia Daniella--Stone'a udowodniono w \atlasref{ME-32::it:thm:daniell-stone}{ME-32: twierdzenie Daniella--Stone'a} przez konstrukcję miary zewnętrznej, mierzalność zbiorów poziomicowych i aproksymację kratową, wykorzystując \atlasref{ME-03::thm:caratheodory}{ME-03: twierdzenie Carathéodory'ego}. Tam też znajduje się \atlasref{ME-32::it:thm:daniell-stone-general}{ME-32: wersja bez jedynki i bez założenia skończoności miary}. Wykaz nie zastępuje szczegółowej kontroli zależności dowodów; migracja nie oznacza zakończenia audytów matematycznych.

